% Limitations of the experimental setup
% Tracking known data and methodology limitations discovered during the
% 2026-05-20 audit (docs/research/audit-findings-2026-05-20.md). Items below
% are documented-as-limitations rather than fixed in this paper's cycle;
% they are queued for follow-up work.

\section{Limitations and Threats to Validity}
\label{sec:limitations}

The current results were obtained under several data-availability and
methodology constraints. We list these explicitly so that downstream readers
can scope conclusions appropriately and so that reproductions can target
the same envelope.

\paragraph{Compound TVL is a placeholder (Finding \#3).}
The Compound V3 RPC fetcher records \texttt{total\_supplied\_usd = 0.0} as a
placeholder pending wiring of the USDC base-units $\to$ USD conversion.
Consequently \texttt{tvl\_compound}, \texttt{debt\_compound}, and the
derived \texttt{dTVL\_compound\_24h} feature column are identically zero
across all 13{,}105 hourly rows of \texttt{joined\_clean.parquet}. Branch~B
of the forecaster therefore receives one constant input
($\approx 14\%$ of its 7-dimensional input capacity), which the model
ignores after the first few epochs but which represents wasted capacity.

\paragraph{Ethereum gas price is a constant stub (Finding \#4).}
\texttt{data/fetch\_gas\_eth.py} is currently a stub
(\texttt{raise NotImplementedError}); the training pipeline fills
\texttt{gas\_gwei = 30.0} as a constant for the entire data window. As with
Compound TVL, this is one further constant Branch~B input; the MCDM
\texttt{f\_cost} factor in the strategy still uses the live spot gas at
decision time, so live-trading behavior is unaffected, but the
backtest cost-normalization is identical across the train/val/test
windows and does not reflect realistic gas variation.

\paragraph{Aggregate-only metrics, no per-quarter breakdown (Finding \#9).}
The headline metrics in \texttt{backtest/run\_main.py} (net APY, annual
Sharpe, max drawdown, Calmar, forecast hit-rate) are computed over the
full test window as a single aggregate. Given the regime structure
documented in \texttt{CLAUDE.md} (Q4 2024 standard deviation 5.75pp vs
2026 Q1 0.57pp --- approximately $10\times$ volatility variation), this
aggregation can mask edge concentration in a single regime. A per-quarter
breakdown table is queued for the next paper cycle.

\paragraph{Composite-loss components logged in unscaled form (Finding \#10).}
\texttt{forecaster/losses.py} logs the bare $\textsc{MSE}$,
$1 - \rho_{w}$, and quantile-loss components to MLflow rather than the
scaled $\alpha \cdot \textsc{MSE}$, $\beta \cdot (1 - \rho_w)$,
$\gamma \cdot Q_{90}$ contributions. This made diagnosing the R\textsuperscript{2}
scale issue (Finding \#1) slower than necessary. Scaled-component logging
is a low-risk addition deferred to a follow-up.

\paragraph{Tie-break bias in winner-series accounting (Finding \#11).}
\texttt{backtest/run\_main.py:\_winner\_series} uses
\texttt{aave >= comp} (rather than strict \texttt{>}), which assigns the
``winner'' label to Aave on ties. In real data ties occur at machine
precision only and the effect on \texttt{n\_rebalances} is bounded by
$\pm 1$; we document this as a convention.

\paragraph{Boundary leakage from 24-hour pct\_change feature (Finding \#12).}
\texttt{data/features.py} computes
\texttt{dTVL\_aave\_24h = tvl\_aave.pct\_change(24)} on the full joined
panel before the train/val/test split is applied, then slices by row
index. The first 24~rows of each non-training fold therefore use TVL
inputs from the previous fold. The contamination is bounded at
$24/2501 \approx 0.96\%$ of test samples; we report it as a
threat-to-validity rather than re-engineer the feature pipeline
per-fold for this cycle.

\paragraph{$N$-way allocation across the lending market.}
Our empirical study covers the two largest Ethereum L1 lending protocols
(Aave V3 + Compound V3, jointly $41\%$ of the $\approx\$54$\,B total
lending TVL --- DeFiLlama April 2026 \cite{defillama2026snapshot}; see
\Cref{tab:lending-tvl}). The dual-branch architecture generalizes to
$N$ protocols without redesign --- each additional protocol requires
only its protocol-specific kink function $f_{\text{kink}}^{(i)}(u)$ and
a per-protocol Branch~B feature wiring.  Four immediate Ethereum L1
targets, ordered by TVL: Spark (SparkLend, \$6.8\,B; Maker-related),
Morpho Blue (\$4.9\,B; permissionless markets), Fluid (\$1.6\,B), and
Euler V2 (\$0.89\,B).  Cross-chain extensions (JustLend on Tron,
\$2.4\,B; Kamino on Solana, \$1.1\,B) defer the rate-fetching
infrastructure work but pose no methodological obstacle.  We leave the
$N$-way extension as the obvious next step.
